\documentclass[12pt, letterpaper]{article}
\date{\today}
\usepackage[margin=1in]{geometry}
\usepackage{amsmath}
\usepackage{hyperref}
\usepackage{cancel}
\usepackage{amssymb}
\usepackage{fancyhdr}
\usepackage{pgfplots}
\usepackage{booktabs}
\usepackage{pifont}
\usepackage{amsthm,latexsym,amsfonts,graphicx,epsfig,comment}
\pgfplotsset{compat=1.16}
\usepackage{xcolor}
\usepackage{tikz}
\usetikzlibrary{shapes.geometric}
\usetikzlibrary{arrows.meta,arrows}
\newcommand{\Z}{\mathbb{Z}}
\newcommand{\N}{\mathbb{N}}
\newcommand{\R}{\mathbb{R}}
\newcommand{\Q}{\mathbb{Q}}
\newcommand{\C}{\mathbb{C}}
\newcommand{\F}{\mathbb{F}}

\newcommand{\Po}{\mathcal{P}}
\newcommand{\Pro}{\mathbb{P}}
\author{Alex Valentino}
\title{CS 336 homework}
\pagestyle{fancy}
\renewcommand{\headrulewidth}{0pt}
\renewcommand{\footrulewidth}{0pt}
\fancyhf{}
\rhead{
	Homework 1\\
	CS 336	
}
\lhead{
	Alex Valentino\\
}
\begin{document}
\begin{itemize}
	\item{Problem 1} Consider the following bank database with 6 tables. For each table, the primary key is underlined.
	\begin{center}
	$\begin{array}{c|c}
	Table & Attributes\\
	\hline
	branch & (\underline{branch\_name}, branch\_city, assets)\\
	\hline	
	customer & (\underline{ID}, customer\_name, street, city)\\
	\hline	
	loan & (\underline{loan\_number}, branch\_name, amount)\\
	\hline
	borrower & (\underline{CustomerID, loan\_number})\\
	\hline
	account & (\underline{account\_number}, branch\_name, balance)\\
	\hline
	depositor & (\underline{CustomerID, account\_number})
	\end{array}	$
	\end{center}
	\begin{enumerate}
		\item \textit{Find the ID and name of customers who live in "Boston".}\\
		$\Pi_{ID, customer\_name}(\sigma_{city = "Boston"}(customer))$
		\item \textit{Find the ID of each borrower who has a loan in branch "Edison."}\\
		$ \Pi_{CustomerID}(\sigma_{branch\_name = "Edison"}(borrower \bowtie_{borrower.loan\_number = loan.loan\_number} loan))$
		\item \textit{Find the ID and names of customers who have deposits in "Edison" branch but do not have loans in "Edison" branch.}\\
		$\Pi_{CustomerID}(\sigma_{branch\_name = "Edison"}) -$\\
		$ \Pi_{ID, customer\_name}(\sigma_{branch\_name = "Edison"}($\\$customer \bowtie_{customer.ID = borrower.CustomerID} borrower \bowtie_{borrower.loan\_number = loan.loan\_number} loan))$
	\end{enumerate}
	\item{Problem 2}
	\begin{enumerate}
		\item 
	\end{enumerate}
\end{itemize}
\end{document}
